\documentclass[12pt,a4paper]{article}

% ---------- پکیج‌ها (همه قبل از xepersian) ----------
\usepackage{fontspec}
\usepackage{amsmath, amssymb, amstext}
\usepackage{graphicx}
\usepackage{geometry}
\usepackage{xcolor}
\usepackage{fancyhdr}

% ---------- تنظیم حاشیه ----------
\geometry{margin=1in}

% ---------- تنظیم fancyhdr ----------
\pagestyle{fancy}
\setlength{\headheight}{15pt} % حل خطای headheight

% هدرها
\lhead{کوییز درس مکانیک‌تحلیلی}
\chead{دانشکده فیزیک دانشگاه صنعتی خواجه نصیرالدین طوسی}
\rhead{\today}

% ---------- بسته فارسی و فونت ----------
\usepackage{xepersian}
\settextfont{B Nazanin}

\begin{document}
	
	% جدول اطلاعات فردی
	\begin{tabular}{|p{7cm}|p{7cm}|}
		\hline
		\textbf{نام و نام خانوادگی:} & \textbf{شماره دانشجویی:} \\
		\hline
	\end{tabular}
	
	% بخش سوالات
	\section*{سوالات:}
	
	یک موشک تک‌مرحله‌ای با مشخصات زیر داریم:  
	\begin{itemize}
		\item جرم اولیه: \(M_0\)
		\item جرم سوخت: \(M_f\)
		\item پیشران ثابت: \(T\)
		\item سرعت خروج گاز: \(u\)
		\item شتاب گرانش: \(g\)
		\item مقاومت هوا: \(F_{\text{\lr{drag}} = k v^2\)
	\end{itemize}
	
	\begin{enumerate}
		\item معادله حرکت موشک را بنویسید، شامل اثر وزن، پیشران و مقاومت هوا.
		\item نشان دهید که بدون مقاومت هوا، تغییر سرعت (\(\Delta v\)) موشک به صورت زیر است:
		\[
		\Delta v = u \ln \frac{M_0}{M_0 - M_f}
		\]
		\item شرط برخاست موشک را بنویسید (حداقل \(T\) لازم برای بلند شدن از زمین).
		\item اگر بخواهیم شتاب اولیه بیش از \(2g\) نشود، شرط مربوط به \(T\) را بنویسید.
	\end{enumerate}
	
	\vspace{10pt}
	\textbf{موفق باشید.}
	
\end{document}
